\documentclass{article}
\usepackage{blindtext}
\usepackage{geometry}
\usepackage{graphicx}
\usepackage[utf8]{inputenc}
\graphicspath{{./figs/}}
\usepackage{multirow, multicol, longtable}
\usepackage{url}
\usepackage[colorlinks=true, linkcolor=blue, citecolor=blue, urlcolor=blue]{hyperref}
\geometry{
	a4paper,
	total={170mm,257mm},
	left=20mm,
	top=20mm,
}
\usepackage[authoryear]{natbib} 
\usepackage{longtable}
\usepackage{caption}
\usepackage{subcaption}
\usepackage{algpseudocode}
\usepackage{algorithm}
\usepackage{anyfontsize}
\usepackage{lineno}
\usepackage{authblk}
\linenumbers
\usepackage{setspace}
\usepackage{booktabs}
\usepackage{tabularx}
\usepackage{amsmath, amsfonts, booktabs, textcomp}
\usepackage{siunitx, microtype}

\title{
	A High-Performance Rule-Based Multi-Agent Framework for Spatio-Temporal 
	Vector Dynamics: Application to \textit{Anopheles stephensi} in Ethiopia}

\author[1,2]{Salvador D. Atagong}
\author[1]{Komi Mensah Agboka}
\author[1]{Kennedy Senagi}
\author[4]{Mark Wamalwa}
\author[2,3]{Henri Tonnang}
\author[2]{John Odindi}

\affil[1]{International Center of Insect Physiology and Ecology (ICIPE)}
\affil[2]{University of KwaZulu-Natal, South Africa}
\affil[3]{International Institute of Tropical Agriculture (IITA), Ibadan, Nigeria}
\affil[4]{University of Nairobi, Kenya}
\date{}

\begin{document}
\maketitle

\begin{abstract}
\textbf{Background:} The invasive malaria vector \textit{Anopheles stephensi} threatens urban Africa, but data scarcity 
limits conventional modelling.  
\textbf{Methods:} We developed a modular agent-based framework coupling a temperature-dependent Metzler matrix life-cycle 
model with spatial environmental layers (population, buildings, climate), behavioural rules, and a lazy spatial grid. 
The model was parameterised for \textit{An. stephensi} and run for six months over two Ethiopian regions (Dire Dawa and Oromia).  
\textbf{Results:} Populations reached near-equilibrium ($\lambda_{\max}=1.003$-$1.074$) with stage durations matching laboratory 
data. Urban cores showed highest suitability (Continuous Boyce Index up to $0.633$). Adult abundance correlated weakly with 
temperature ($r=-0.12$, $p=0.12$) but significantly with precipitation in Dire Dawa ($r=-0.72$, $p<0.001$).  
\textbf{Conclusion:} The framework reproduces known \textit{An. stephensi} ecology and identifies urban hotspots. Its 
YAML-driven (YAML is a human-readable, lightweight data serialization language used primarily for configuration files, 
data exchange, and DevOps tools like Kubernetes and Ansible) parameterisation makes it adaptable to other species and regions, 
supporting evidence-based vector control in data-limited settings.

\paragraph{Keywords:} multi-agent systems, vector-borne disease, 
\textit{Anopheles stephensi}, Metzler matrix, high-performance computing, spatial ecology.
\end{abstract}

\section{Introduction}
\label{sec:introduction}
\par
Vector-borne diseases (VBDs) represent one of the most profound challenges to global health security, accounting for 
approximately 700,000 deaths annually and imposing significant economic burdens on developing nations \citep{2025.RL.HMD}. 
Malaria alone threatens nearly half the world's population, with Sub-Saharan Africa bearing 94\% of the global case burden 
as \citep{2024.WH.WWM}. While the past two decades witnessed substantial reductions in malaria mortality through coordinated 
interventions, this hard-won progress faces unprecedented threats from ecological and demographic shifts. 
The rapid urbanization of Sub-Saharan Africa, currently the world's fastest urbanizing region \citep{2019.IS.UAI}, 
has created novel transmission ecologies that challenges historical assumptions about vector ecology and control 
strategies \citep{2024.NK.MBD}.
\par
The primary driver of this epidemiological transition is the globalization of previously localized vectors, particularly 
the invasive malaria mosquito \textit{Anopheles stephensi}. Unlike endemic African malaria vectors that predominantly 
breed in natural water bodies, \textit{An. stephensi} exhibits remarkable ecological plasticity, thriving in artificial 
containers and man-made structures characteristic of urban environments \citep{2025.GC.SAB}. This species shares critical 
adaptive traits with \textit{Aedes aegypti}, the primary dengue vector, including container breeding, anthropophily, and 
exophilic resting behaviors \citep{2022.AH.TPI}. The consequences of this adaptation were first recorded in Djibouti, 
where\textit{ An. stephensi} establishment triggered a 30-fold increase in malaria cases within seven years, transforming a 
nation on the verge of malaria elimination into a high-transmission zone \citep{2022.AH.TPI}. This vector's capacity to exploit 
previously low-risk urban settings necessitates fundamentally new approaches to surveillance, modeling, and 
intervention design \citep{2024.TM.AQA}.
\par
A critical bottleneck in responding to this emerging threat is the pervasive "Data Gap" characterizing many resource-constrained 
settings where \textit{An. stephensi} is establishing. Effective vector control relies on robust decision-support systems built 
upon high-quality, high-resolution entomological and epidemiological data, yet national surveillance programs often operate with 
limited financial and technological resources \citep{2025.RL.HMD}. Current modeling approaches struggle to bridge the gap between 
theoretical elegance and practical utility: empirical models propagate measurement errors through unreliable data streams, while 
process-based models frequently lack the spatial granularity and computational efficiency needed for operational 
deployment \citep{2025.RL.HMD}. This disconnect is particularly acute for invasive vectors, where absence of historical 
baseline data undermines traditional statistical approaches \citep{2019.EL.ALO}.
\par
To address these challenges, we propose a novel computational framework that integrates high-performance computing methodologies 
with ecological process modeling. Our system implements a layered and modular, rule-based Multi-Agent Simulation (MAS) architecture that 
harmonizes heterogeneous spatio-temporal data streams into a unified analytical environment. This framework is part of the 
family of ABS (Agent Based Simulation) which are simulation models where the researcher explicitly describe the decision 
processes of simulated actors at micro-level \citep{2008.PE.ITM}. The proposed framework bridges Computer Science and 
Environmental Science by conceptualizing environmental data as dynamic information layers, including static rasters (for instance: 
elevation, building density, population density), time-series climate variables, and vector data (shapefiles), that interact with 
autonomous agent populations. We distinguish between two agent classes: "Inert Agents" representing static environmental features 
(urban water containers serving as breeding sites) and "Living Agents" modeling biological vectors (\textit{An. stephensi} mosquitoes for 
example) with behavioral rules derived from empirical studies.
\par
This study demonstrates how such synthetic modeling approaches can provide critical insights where field data are scarce or 
unreliable. Using Dire Dawa, and Oromia woredas from Ethiopia, which are regions experiencing \textit{An. stephensi} invasion, 
\citep{2024.GZ.ROC,2023.DH.FRO} as a proof-of-concept, we simulate 6 months of mosquito population dynamics across the whole regions. 
Our results illustrate how agent-based modeling can identify high-risk micro-habitats, simulate seasonal population trajectories, and 
potentially assist intervention strategies in data-limited settings. By making the simulation framework open-source and computationally 
efficient, we aim to empower public health agencies in resource-constrained regions with decision-support tools capable of addressing the 
unique challenges posed by urban-adapted invasive vectors.


\section{Materials and Methods}
\label{sec:methods}
\subsection{Framework Architecture}
\label{sec:framework}
\subsubsection{Overview}
% comprehensive chart of the framework
The simulation framework follows a modular, agent-based design that integrates spatially explicit environmental data, 
temperature-dependent life-cycle models, and behavioural rules. It is implemented in Java and built around four core 
subsystems: (i) a time manager that advances simulation time in discrete ticks (15-minutes per tick) and interpolates 
climate data; (ii) a spatial registry that indexes agents on a lazy grid for efficient neighbourhood queries; 
(iii) an agent layer system that manages agents and evaluates rules in parallel; and (iv) a lifecycle model that computes 
temperature-driven development and survival rates. The overall architecture is illustrated in Figure \ref{fig:arch}.
\begin{figure}[!htb]
\centering
\includegraphics[width=\textwidth]{figs/phd_graphs-rmas.drawio.png}
\caption{High-level architecture of the simulation framework. The diagram is organised into three layers: 
\textit{Data Layer} (input data sources), \textit{Core Engine Layer} (simulation components), and \textit{Analysis 
\& Visualization Layer} (post-processing and validation). Solid arrows denote \textit{data flow} 
(e.g., configuration files, rasters, climate NetCDF, and simulation outputs). Dashed arrows denote 
\textit{control flow} (e.g., tick scheduling, spatial queries, rule evaluation, and lifecycle updates). The analysis 
layer reads CSV snapshots generated by the persistence service and produces daily summaries, suitability maps, 
validation metrics, and final figures.}
\label{fig:arch}
\end{figure}

\subsection{Core Components}
\label{sec:core_components}

\subsubsection{Agent layers}
Two agent types are implemented: \texttt{LivingAgent} (representing a mosquito) and \texttt{InertAgent} (representing a 
breeding site). Living agents track life stage (egg, larva, pupa, adult), age, energy, gravid status, and resting behaviour.
Their movement is modelled as a random walk with a step size derived from field observations of flight range \citep{1966.MA.FRL}. 
Inert agents maintain water volume, egg count, larval count, and carrying capacity; they update their capacity dynamically based on 
water level, mimicking density-dependent mortality observed in larval habitats \citep{2016.ST.OTI}. Agent layers are processed in 
parallel using a partition-based container (\texttt{AgentContainer}) and a rule engine that evaluates behavioural rules (e.g., 
feeding, oviposition, resting) in priority order \citep{2023.ST.ASR}.

\subsubsection{Raster layers and climate data}
Static raster layers (elevation, building density, population density) are loaded from GeoTIFF files using GeoTools and stored either 
in memory or memory-mapped for large grids. Climate time series (2-m temperature and total precipitation) are 
ingested from NetCDF files (ERA5-Land reanalysis) into interpolated raster layers that perform linear temporal interpolation between 
hourly steps \citep{2014.GC.AED, 2017.CL.IMS}.

\subsubsection{Lifecycle model}
\label{sec:lifecycle}
The lifecycle model represents mosquito population dynamics using a discrete-time, stage-structured Metzler matrix 
\citep{2023.ST.ASR,2025.KM.WTR}. Four developmental stages are tracked: eggs ($E$), larvae ($L$), pupae ($P$), and adults ($A$). 
All rates are temperature-dependent and defined on a per-day basis, then converted to tick-level probabilities via the Poisson 
process:
\[
p = 1 - \exp(-r \Delta t),
\]
where $\Delta t$ is the tick duration in days ($\Delta t = 0.25$ h $= 0.01042$ days). 
The daily projection matrix is:
\[
\mathbf{M}(T) = 
\begin{pmatrix}
(1 - d_E) S_E & 0 & 0 & F \cdot K \\
d_E S_E & (1 - d_L) S_L & 0 & 0 \\
0 & d_L S_L & (1 - d_P) S_P & 0 \\
0 & 0 & d_P S_P & 1 - \mu_A
\end{pmatrix},
\]
with stage abundances evolving as $[E,L,P,A]_{t+1}^\top = \mathbf{M}(T) [E,L,P,A]_t^\top$. The temperature-dependent functions are:
\begin{itemize}
  \item \textbf{Development rates} (day$^{-1}$): $d_X(T) = \max(0, a_X T^2 + b_X T + c_X)$ for $X \in \{E,L,P\}$.
  \item \textbf{Stage survival} (Gaussian): $S_X(T) = \text{amp}_X \cdot \exp\left(-\frac{1}{2}\left(\frac{T - \mu_X}{\sigma_X}\right)^2\right)$.
  \item \textbf{Adult mortality} (day$^{-1}$): $\mu_A(T) = \max(0, a_A T^2 + b_A T + c_A)$.
  \item \textbf{Fecundity} (eggs/female/day): $F(T) = \max\left(0, a_F e^{b_F T} - e^{b_F T_{\max} - \left(\frac{T_{\max}-T}{c_F}\right)^2}\right)$.
  \item \textbf{Host carrying capacity}: $K = 1 + \ln(1 + \text{population density})$, which scales fecundity based on human population 
  density (diminishing returns at high density).
  \item \textbf{Density-dependent larval mortality}: each aquatic habitat has a carrying capacity $\text{cap}$ (linearly scaled with water 
  volume); if larval count exceeds $\text{cap}$, excess larvae are removed probabilistically.
\end{itemize}
All temperature-response coefficients are stored in a YAML configuration file under the \texttt{species} key, making the model generic 
and easily adaptable to any mosquito species by replacing the coefficients in Table \ref{tab:lifecycle_params}. This parametrisation 
follows the principles of generic insect population modelling \citep{2025.KM.DPF,2025.KM.WTR,2026.KM.THM}.

\begin{table}[!htb]
\centering
\caption{Generic temperature-response parameters for a mosquito lifecycle model.}
\label{tab:lifecycle_params}
\begin{tabular}{llll}
\toprule
Parameter set & Meaning & Unit & YAML key \\ \midrule
$d_E$ coeff. & Egg development rate & $1/^\circ\text{C}^2, 1/^\circ\text{C}, 1$ & \texttt{egg\_dev\_a,b,c} \\
$d_L$ coeff. & Larva development rate & $1/^\circ\text{C}^2, 1/^\circ\text{C}, 1$ & \texttt{larva\_dev\_a,b,c} \\
$d_P$ coeff. & Pupa development rate & $1/^\circ\text{C}^2, 1/^\circ\text{C}, 1$ & \texttt{pupa\_dev\_a,b,c} \\
$S_E$ & Egg survival (amp, mean, sigma) & dim., $^\circ\text{C}, ^\circ\text{C}$ & \texttt{egg\_survival\_...} \\
$S_L$ & Larva survival (amp, mean, sigma) & dim., $^\circ\text{C}, ^\circ\text{C}$ & \texttt{larva\_survival\_...} \\
$S_P$ & Pupa survival (amp, mean, sigma) & dim., $^\circ\text{C}, ^\circ\text{C}$ & \texttt{pupa\_survival\_...} \\
$\mu_A$ coeff. & Adult mortality & $1/^\circ\text{C}^2, 1/^\circ\text{C}, 1$ & \texttt{adult\_mort\_a,b,c} \\
$F$ coeff. & Fecundity & $a_F, b_F, T_{\max}, c_F$ & \texttt{fecundity\_a,b,Tmax,c} \\
$K$ & Host scaling factor & dimensionless & (derived) \\ \bottomrule
\multicolumn{4}{p{0.9\linewidth}}
{\footnotesize \textit{Note:} 
"coeff." denotes the three coefficients $(a,b,c)$ of the quadratic function. 
"1" indicates dimensionless (unit 1). 
"dim." stands for dimensionless. 
For survival parameters, "amp" is dimensionless, "mean" and "sigma" are in \si{\celsius}.
} \\
\end{tabular}
\end{table}

\subsubsection{Rule engine and spatial registry}
Behavioural rules (e.g., "\texttt{feed}", "\texttt{lay\_eggs}", "\texttt{rest\_in\_building}") are defined in a YAML configuration 
file and evaluated using a JavaScript engine (GraalVM \citep{2023.VT.CPA}). The spatial registry implements a lazy grid with a configurable cell size 
(default $0.001^\circ \approx 110$ m). Neighbour queries are performed by checking cells within a Chebyshev distance 
$\lceil \text{radius}/\text{cellSize}\rceil$; results are cached for one second per cell to reduce computational cost.

\subsubsection{Simulation Configuration}
% how to configure a simulation
\par
All simulation parameters are defined in a single YAML file (\texttt{simulation.yaml}). The file specifies:  
\par
\begin{itemize}
  \item \textbf{Time settings}: start date, total number of ticks, tick duration (15-minutes).  
  \item \textbf{File paths}: study area shapefile, elevation, building density, population rasters, and the climate NetCDF.  
  \item \textbf{Species parameters}: coefficients for development, survival, fecundity, and adult mortality.  
  \item \textbf{Tokens}: mapping between environment variable names (e.g., "temperature") and raster layers.  
  \item \textbf{Agent layers and rules}: each layer contains a list of rules with a condition, an action, and a priority.  
  \item \textbf{Seeding}: number of breeding sites and mosquitoes, habitat grid resolution, and whether to seed uniformly or 
  using a weighted habitat suitability map.  
  \item \textbf{Project defaults}: agent search radius, step size, maximum age.
\end{itemize}
\par
At startup, the simulator loads the configuration, initialises the time manager, loads rasters and climate data, 
creates the spatial registry, builds agent layers, seeds the initial population, and starts the simulation engine.

\subsubsection{Result Analysis}
\label{sec:analysis}
% details on the analysis notebook how are simulation result analyzed details
Simulation outputs are exported as CSV snapshots every 100 ticks. Each snapshot contains per-agent records 
with tick count, agent ID, layer, type, coordinates, alive flag, age, stage, energy, gravid status, and for inert agents, egg/larval 
counts and water volume. Environmental values (temperature, precipitation, etc) are attached to each agent's row via bilinear interpolation 
from the current climate frame.
\par
The analysis, visulization and validation of the resutls are realized through a jupyter notebook (Python) which performs the following steps:
\par
\begin{enumerate}
  \item \textbf{Daily aggregation}: The CSV is read in chunks (100 000 rows) to compute daily adult, larva, and pupa counts, as well 
  as breading sites statistics (mean water volume, total eggs/larvae), using the tick-to-day conversion ($96$ ticks per day by default).  
  \item \textbf{Population dynamics}: Stage counts are plotted, and stage durations are estimated from cross-correlation of smoothed 
  population peaks. Theoretical durations are computed from species parameters at the mean observed temperature.  
  \item \textbf{Environmental correlation}: Pearson correlation between daily adult counts and mean temperature / precipitation is calculated.  
  \item \textbf{Suitability mapping}: A kernel density estimate (KDE) of adult points (bandwidth 0.002° ≈ 220 m) is generated and exported 
  as a GeoTIFF. A breeding site potential raster is created by rasterising the sum of eggs and larvae in water tanks on the last tick.  
  \item \textbf{Validation}: The continuous Boyce index \citep{2006.HA.ETA} and the proportion of occurrence points falling into the top 20\% 
  of the suitability surface are computed using independently collected occurrence data (from GBIF, iNaturalist, literature, clipped to 
  the study area).  
  \item \textbf{Export}: Daily summaries and a sample of raw data are saved for further analysis.
\end{enumerate}
\par
All figures and rasters are stored in a user-defined output directory.

\subsection{Case Study Design}
\label{sec:case_study}
\subsubsection{Vector and Study Site}
\par
The case study is set in Ethiopia, focusing on two contrasting regions: Dire Dawa and Oromia. 
Both areas have experienced recent incursions of the Asian malaria vector \textit{Anopheles stephensi} and are 
considered high-risk zones for urban and peri-urban malaria transmission \citep{2020.MB.GDO, 2021.MB.AUO, 2023.DH.FRO}.
\par
Dire Dawa (\ang{9.6049} N, \ang{41.8585} E) is a low-altitude (\SI{\sim 1200}{\meter}) urban centre with a semi-arid climate and a known 
history of \textit{An. stephensi} establishment, first reported in 2018 \citep{2020.MB.GDO}. The Oromia (\ang{7.5460} N, \ang{40.6347} E) 
study area comprises several \textit{woredas} (e.g., Adama, Bishoftu, and surrounding districts) that span an altitudinal gradient 
from \SIrange{1500}{2500}{\meter}, representing a mix of agricultural, peri-urban, and highland ecosystems where the vector has been detected since 
2019 \citep{2021.MB.AUO, 2024.TA.SDA}. These sites were selected because they provide a natural experiment to evaluate the simulator's 
ability to capture spatial heterogeneity in vector abundance driven by temperature, urbanisation, and host availability 
\citep{2025.TD.BOA, 2024.GZ.ROC}.

\subsubsection{Vector Biology}

\textit{Anopheles stephensi} is a highly adaptive urban malaria vector native to South Asia and the Middle East, now established 
in the Horn of Africa \citep{2024.RT.IAS}. Its life cycle is tightly coupled to temperature, and its larvae exploit a wide range 
of artificial containers, especially overhead water tanks and construction pits \citep{2016.ST.OTI}. 

The temperature-dependent functions described in Section \ref{sec:lifecycle} were parameterised for \textit{An. stephensi} using 
laboratory and field data from the literature (Table \ref{tab:stephensi_params}). Development rates ($d_E, d_L, d_P$) and adult 
mortality ($\mu_A$) were fitted to the observations of \citet{2023.JJ.BBA} and \citet{2025.CG.EOT}. Survival probabilities 
($S_E, S_L, S_P$) were derived from the meta-analysis of \citet{2022.OC.TIT}. Fecundity coefficients ($a_F, b_F, T_{\max}, c_F$) 
were taken from \citet{2023.JJ.BBA}, who estimated daily egg production under varying thermal regimes. 

The host carrying capacity $K$ uses human population density (obtained from the WorldPop raster) as a proxy for biting resources, 
assuming that \textit{An. stephensi} is highly anthropophilic in urban settings \citep{2024.TA.SDA, 2022.AH.TPI}. Density-dependent 
larval mortality follows the tank capacity model of \citet{2016.ST.OTI}, where capacity is linearly scaled with water volume 
(maximum 500 larvae at \SI{100}{\percent} water volume).

The resulting daily projection matrix for \textit{An. stephensi} at a reference temperature $T = \SI{25}{\celsius}$ is:
\[
\mathbf{M}(25) = \begin{pmatrix}
0.48 & 0 & 0 & 8.2 \\
0.52 & 0.76 & 0 & 0 \\
0 & 0.24 & 0.70 & 0 \\
0 & 0 & 0.30 & 0.91
\end{pmatrix}
\]
where the dominant eigenvalue $\lambda_{\max} = 1.03$, indicating slow population growth under optimal conditions, consistent with 
field observations in Ethiopia \citep{2021.MB.AUO, 2022.AH.TPI}.

\begin{table}[!htb]
\centering
\small % Slighting smaller font to ensure table fits page width
\caption{Species-specific parameters for \textit{Anopheles stephensi} used in the case study.}
\label{tab:stephensi_params}
\begin{tabular}{@{}llll@{}}
\toprule
Parameter & Meaning & Value / coefficients & Source \\ \midrule
$d_E$ & Egg development (day$^{-1}$) & $a=-0.00098, b=0.0502, c=-0.348$ & \citet{2022.OC.TIT} \\
$d_L$ & Larva development (day$^{-1}$) & $a=-0.00082, b=0.0425, c=-0.305$ & \citet{2023.JJ.BBA} \\
$d_P$ & Pupa development (day$^{-1}$) & $a=-0.00074, b=0.0398, c=-0.291$ & \citet{2022.OC.TIT} \\
$S_E$ & Egg survival & $\text{amp}=0.94, \mu=26.2, \sigma=3.5$ & \citet{2022.OC.TIT} \\
$S_L$ & Larva survival & $\text{amp}=0.88, \mu=26.8, \sigma=3.8$ & \citet{2022.OC.TIT} \\
$S_P$ & Pupa survival & $\text{amp}=0.91, \mu=26.5, \sigma=3.6$ & \citet{2022.OC.TIT} \\
$\mu_A$ & Adult mortality (day$^{-1}$) & $a=0.00012, b=-0.0018, c=0.055$ & \citet{2025.CG.EOT} \\
$F$ & Fecundity (eggs/female/day) & $a_F=-0.0099, b_F=0.567, T_{\max}=32.0, c_F=2.97$ & \citet{2023.JJ.BBA} \\
$K$ & Host scaling factor & $1 + \ln(1+\text{popdens})$ & - \\ \bottomrule %Worldpop(2020)
\end{tabular}
\end{table}


\subsubsection{Data Sources}
% Required packages: \usepackage{siunitx, microtype, booktabs}

\label{sec:data_sources}
The simulation integrates five categories of data:

\begin{itemize}
    \item \textbf{Static environmental rasters} -- Elevation (SRTM, \SI{30}{\meter}), building density (Global 
	Human Settlement Layer), and human population density (WorldPop, 2020) were resampled to a common \ang{0.001} 
	(\SI{\sim 110}{\meter}) grid over the study area.
    
    \item \textbf{Climate time series} -- Hourly \SI{2}{\meter} temperature and total precipitation for 2020 were 
	obtained from the ERA5-Land reanalysis (Copernicus Climate Change Service) and stored as NetCDF files.
    
    \item \textbf{Vector occurrence data} -- Presence points for \textit{An. stephensi} in Ethiopia were compiled from 
	the Vector Atlas database \citep{2025.VE.VAP}, GBIF \citep{2011.AN.BIG}, and recent entomological surveys 
	\citep{2020.MB.GDO, 2021.MB.AUO, 2023.DH.FRO}. Only records with precise coordinates and collection dates 
	in 2018--2022 were retained.
    
    \item \textbf{Bionomic and life-history data} -- Parameter values (Table \ref{tab:stephensi_params}) were extracted 
	from peer-reviewed literature \citep{2022.OC.TIT, 2023.JJ.BBA, 2025.CG.EOT} and supplemented with expert knowledge.
    
    % \item \textbf{Hydrographic information} -- Locations of potential breeding sites (overhead tanks, wells, construction 
	% pits) were derived from high-resolution satellite imagery and local surveys \citep{2016.ST.OTI, 2024.TA.SDA}.
\end{itemize}

All data were harmonised to the WGS84 geographic coordinate system (EPSG:4326) and stored in a structured file hierarchy.

\subsubsection{Simulation Assupmtions, Rules and Configuration}
\label{sec:sim_assumptions}

The simulation makes the following key assumptions:
\begin{enumerate}  
    \item \textbf{Temperature is the primary abiotic driver} -- All development, survival, and reproduction rates respond 
	only to ambient temperature; effects of relative humidity, rainfall, and other variables are omitted or implicitly 
	captured through breeding site water dynamics.  
    
    \item \textbf{Host availability scales with human population density} -- The effective fecundity $F \cdot K$ uses 
	population density as a proxy for blood-meal availability, assuming that females feed preferentially on humans in 
	urban settings \citep{2022.AH.TPI}.  
    
    \item \textbf{Aquatic habitats are static} -- Water tank positions and their maximum water volume are fixed at 
	initialisation; only evaporation and rainfall modify water level over time.  
    
    \item \textbf{Mosquito movement is localised} -- Adults disperse via a random walk with a fixed step size 
	(\ang{0.00005} $\approx$ \SI{5.5}{\meter} per tick) and a search radius of \ang{0.005} (\SI{\sim 550}{\meter}), 
	reflecting limited flight range \citep{1966.MA.FRL}.  
    
    \item \textbf{Resting and feeding follow simple rules} -- Adults rest indoors or outdoors based on building 
	density and temperature, and they feed when energy falls below a threshold, with success depending on 
	local population density.
\end{enumerate}

The behavioural rules (Table \ref{tab:rules}) are encoded in the YAML configuration file (\texttt{simulation.yaml}) and 
evaluated in priority order at each tick. The complete configuration used for the case study is provided as supplementary 
material.

\begin{table}[!htb]
\centering
\caption{Behavioural rules for \textit{An. stephensi} adults.}
\label{tab:rules}
\begin{tabular}{@{}llc@{}}
\toprule
Condition & Action & Priority \\ \midrule
Stage = Adult \textbf{and} energy $<$ 0.6 \textbf{and} not resting & \texttt{feed} & 8 \\
Stage = Adult \textbf{and} energy $>$ 0.7 \textbf{and} not gravid & \texttt{get\_gravid} & 7 \\
Gravid = True \textbf{and} $T > \SI{15}{\celsius}$ & \texttt{lay\_eggs} & 6 \\
$T < \SI{10}{\celsius}$ \textbf{or} $T > \SI{40}{\celsius}$ & \texttt{die} & 9 \\
Precipitation $>$ \SI{0.008}{\meter} \textbf{and} build. dens. $>$ 0.1 \textbf{and} not resting & \texttt{rest\_in\_building} & 5 \\
Stage = Adult \textbf{and} not resting & \texttt{move\_random} & 3 \\
Otherwise & \texttt{rest} & 2 \\ \bottomrule
\end{tabular}
\end{table}

The simulation was run for 180 days (six months, tick = \SI{15}{\minute}, \num{17280} ticks) with a burn-in of 5 days 
to allow transients to decay. Output snapshots were exported every 100 ticks and analysed using the Python notebook described in 
Section \ref{sec:analysis}.

\subsubsection{Experimental Setup}
All simulations were executed on a laptop equipped with an Intel Core i5-6300U processor (2 physical cores, 4 threads, base frequency 
\SI{2.40}{\giga\hertz}, max turbo \SI{3.00}{\giga\hertz}) and \SI{8}{\giga\byte} of RAM. The operating system was Linux (x86\_64), 
and the Java Virtual Machine (JVM) was configured with a maximum heap size of \SI{1950}{\mega\byte} (approximately \SI{80}{\percent} of physical memory). 
The simulation engine used a thread pool with three worker threads (one less than the total number of logical threads) to avoid 
oversubscription. The spatial registry implemented a lazy grid with a cell size of \SI{0.001}{\degree} ($\approx$\SI{110}{\meter}). 
Each simulation ran for 180 simulated days (\num{17280} ticks of \SI{15}{\minute}). Output snapshots were exported every 100 ticks. 
The analysis notebook (Python 3.10) was executed on the same hardware.


\section{Results}
\label{sec:results}

The simulation was executed for a duration of six months (180 days, comprising \num{17280} ticks of \SI{15}{\minute} each) 
over the contrasting study areas of Dire Dawa and  Oromia in Ethiopia. The model utilised the 
temperature-dependent lifecycle parameters specifically fitted for \textit{Anopheles stephensi}. Output snapshots were processed 
to evaluate population stability, environmental response, and spatial suitability.

The mosquito population reached a near-steady state after approximately 10 days of simulation. Across both sites, adult 
abundance stabilised at approximately \num{85000} individuals, while larval and pupal counts fluctuated around \num{35000} and \num{21000}, 
respectively (Figure \ref{fig:population_dynamics}). Analysis of the six-month time series revealed that the daily adult count maintained 
a weak negative correlation with temperature (Pearson's $r = -0.119$, $p = 0.12$) and negligible correlation with precipitation 
($r = 0.009$, $p = 0.91$). These results suggest that under the simulated conditions, intrinsic density-dependent regulation and the 
availability of stable aquatic habitats (water tanks) dominated over short-term climatic fluctuations.

\begin{figure}[!htb]
     \centering
     \begin{subfigure}[b]{0.49\textwidth}
         \centering
         \includegraphics[width=\textwidth]{figs/dire_dawa/mosquito_dynamics.png}
         \caption{Dire Dawa}
     \end{subfigure}
     \hfill
     \begin{subfigure}[b]{0.49\textwidth}
         \centering
         \includegraphics[width=\textwidth]{figs/oromia/mosquito_dynamics.png}
         \caption{Oromia}
     \end{subfigure}
     \caption{Daily counts of adult, larval, and pupal \textit{An. stephensi} over the six-month simulation. Population levels show
	  high stability following the initial 10-day burn-in period.}
     \label{fig:population_dynamics}
\end{figure}

Theoretical stage durations were calculated based on the mean observed temperature (\SI{19.9}{\celsius}), yielding 5.6 days 
for larvae and 8.4 days for pupae. The dominant eigenvalue ($\lambda_{\max}$) of the Metzler projection matrix, averaged over 
the full simulation period, was \num{1.003}. This value indicates a slowly growing population near equilibrium, which is consistent 
with the established status of \textit{An. stephensi} in urban Ethiopia \citep{2021.MB.AUO}. Figure \ref{fig:growth} illustrates the 
moderate agreement between the observed daily adult growth factor and the predicted $\lambda_{\max}$.

\begin{figure}[!htb]
     \centering
     \begin{subfigure}[b]{0.49\textwidth}
         \centering
         \includegraphics[width=\textwidth]{figs/dire_dawa/growth_validation.png}
         \caption{Dire Dawa}
     \end{subfigure}
     \hfill
     \begin{subfigure}[b]{0.49\textwidth}
         \centering
         \includegraphics[width=\textwidth]{figs/oromia/growth_validation.png}
         \caption{Oromia}
     \end{subfigure}
     \caption{Observed daily adult growth factor versus predicted $\lambda_{\max}$ (red dashed line). The average predicted 
	 growth factor ($\lambda_{\max} = 1.003$) aligns with the observed stability.}
     \label{fig:growth}
\end{figure}

A total of \num{1336} water tanks were simulated. Mean water volume remained highly stable at $\SI{85.3 \pm 0.02}{\percent}$, 
as evaporation losses were balanced by periodic precipitation events (Figure \ref{fig:tank_stats}). The reciprocal trend 
observed between egg and larval counts per tank confirms active hatching and stage transition within these permanent aquatic 
habitats.

\begin{figure}[!htb]
     \centering
     \begin{subfigure}[b]{0.49\textwidth}
         \centering
         \includegraphics[width=\textwidth]{figs/dire_dawa/tank_stats.png}
         \caption{Dire Dawa}
     \end{subfigure}
     \hfill
     \begin{subfigure}[b]{0.49\textwidth}
         \centering
         \includegraphics[width=\textwidth]{figs/oromia/tank_stats.png}
         \caption{Oromia}
     \end{subfigure}
     \caption{Water tank statistics: (top) mean water volume; (middle) mean eggs per tank; (bottom) mean larvae per tank.}
     \label{fig:tank_stats}
\end{figure}

Daily average temperatures ranged between \SIrange{19.5}{20.1}{\celsius}, while precipitation remained relatively constant 
at \SI{\sim 1590}{\milli\meter\per\day} (Figure \ref{fig:environment}). The low variability in these drivers over the 
six-month period explains the dominance of endogenous population processes over climatic forcing in this specific case study.

\begin{figure}[!htb]
     \centering
     \begin{subfigure}[b]{0.49\textwidth}
         \centering
         \includegraphics[width=\textwidth]{figs/dire_dawa/environment.png}
         \caption{Dire Dawa}
     \end{subfigure}
     \hfill
     \begin{subfigure}[b]{0.49\textwidth}
         \centering
         \includegraphics[width=\textwidth]{figs/oromia/environment.png}
         \caption{Oromia}
     \end{subfigure}
     \caption{Daily average environmental conditions (Temperature in red, Precipitation in blue) throughout the simulation.}
     \label{fig:environment}
\end{figure}

Spatial suitability maps were generated using Kernel Density Estimation (KDE) with a bandwidth of \ang{0.002} 
(\SI{\sim 220}{\meter}). As shown in Figure \ref{fig:adult_suitability}, adult mosquito density was highest in urban 
cores and areas with significant building density. Validation against independent occurrence points yielded a 
continuous Boyce index of \num{0.633}, confirming the model's moderate predictive power for invasive vector establishment. 
The final breeding site potential, derived from the distribution of eggs and larvae in water tanks, is presented in 
Figure \ref{fig:breeding_potential}.

\begin{figure}[!htb]
     \centering
     \begin{subfigure}[b]{0.49\textwidth}
         \centering
         \includegraphics[width=\textwidth]{figs/dire_dawa/adult_suitability.png}
         \caption{Dire Dawa}
     \end{subfigure}
     \hfill
     \begin{subfigure}[b]{0.49\textwidth}
         \centering
         \includegraphics[width=\textwidth]{figs/oromia/adult_suitability.png}
         \caption{Oromia}
     \end{subfigure}
     \caption{Adult mosquito suitability maps (KDE) overlaid with occurrence records (red dots). High density 
	 correlates with urban infrastructure.}
     \label{fig:adult_suitability}
\end{figure}

\begin{figure}[!htb]
     \centering
     \begin{subfigure}[b]{0.49\textwidth}
         \centering
         \includegraphics[width=\textwidth]{figs/dire_dawa/breeding_site_potential.png}
         \caption{Dire Dawa}
     \end{subfigure}
     \hfill
     \begin{subfigure}[b]{0.49\textwidth}
         \centering
         \includegraphics[width=\textwidth]{figs/oromia/breeding_site_potential.png}
         \caption{Oromia}
     \end{subfigure}
     \caption{Breeding site potential rasterized from final egg and larval counts in water tanks, 
     overlaid with occurrence records (red dots).}
     \label{fig:breeding_potential}
\end{figure}

\newpage
\par
Table \ref{tab:performance} summarizes the computational and agent statistics of the simulation for both study sites.
Both simulations completed the full 180-day period (\num{17280} ticks) without errors. The Oromia run had a slightly higher 
average tick time (\SI{124.9}{\milli\second} vs.\ \SI{116.0}{\milli\second}) and lower effective frame rate (\num{8.0} FPS), 
despite having fewer agents (\num{109436} vs.\ \num{176350}). This is attributable to the larger spatial extent of the Oromia 
study area, which required a much larger spatial grid (\num{60898750} cells vs.\ \num{193410} cells for Dire Dawa), leading to 
a lower average agent density per cell (\num{0.003} vs.\ \num{2.37}). The spatial registry's lazy grid implementation helped 
to maintain acceptable performance even with millions of empty cells.
\par
Memory consumption was higher for the Oromia simulation (\SI{82.8}{\percent} of the \SI{1950}{\mega\byte} heap) due to the 
larger grid metadata and a greater number of rule evaluations (\num{12000} vs.\ \num{7012}). However, both runs stayed well 
below the memory limit, confirming that a standard laptop can execute the model for region-scale simulations. The total number 
of spatial inserts and updates reflects the differences in initial seeding and movement activity between the two sites.
\par
Overall, the simulation framework demonstrated stable performance, with average tick times below \SI{125}{\milli\second}, 
effectively enabling the simulation of six months of vector population dynamics in less than one hour of wall-clock time 
on consumer hardware.

\begin{table}[!htb]
\centering
\caption{Performance metrics and agent counts for the two simulation runs.}
\label{tab:performance}
\begin{tabular}{l S[table-format=8.1] S[table-format=8.1]}
\toprule
Metric & {Dire Dawa} & {Oromia} \\ \midrule
Total runtime (s) & 2499.6 & 2154.9 \\
Average tick time (ms) & 124.9 & 116.0 \\
Effective FPS & 8.0 & 8.6 \\
Peak memory usage (MB) & 910 & 1614 \\
Memory utilization (\% of max) & 46.7 & 82.8 \\
Total agents (end of simulation) & 176350 & 109436 \\
\quad Mosquitoes & 171608 & 108100 \\
\quad Water tanks & 4742 & 1336 \\
Rule evaluations & 7012 & 12000 \\
Actions executed & 19463668 & 21240236 \\
Spatial registry -- total inserts & 505987 & 213813 \\
Spatial registry -- total updates & 251 & 227 \\
Maximum agents per spatial cell & 1416 & 237 \\ \bottomrule
\end{tabular}
\end{table}



\section{Discussion}
\label{sec:discussion}
\par
The simulation results demonstrate that the agent-based model successfully reproduces population dynamics for \textit{Anopheles stephensi} 
across two distinct Ethiopian ecological settings. In the Oromia \textit{woredas}, the population maintained a near-equilibrium state 
($\lambda_{\max} = 1.003$), whereas in Dire Dawa, the model predicted a more robust growth potential ($\lambda_{\max} = 1.074$). These 
findings align with field reports that identify Dire Dawa as a primary hotspot for the species' establishment in urban Ethiopia \citep{2021.MB.AUO}. 
The theoretical stage durations—ranging from \SI{5.0}{\day} in the warmer Dire Dawa (\SI{21.9}{\celsius}) to \SI{5.6}{\day} in Oromia 
(\SI{19.9}{\celsius}) for larvae—fall within the ranges reported in laboratory studies \citep{2023.JJ.BBA}, supporting the validity of our 
temperature-dependent parameterisation.
\par
A key finding is the site-specific response to environmental drivers. In Oromia, adult abundance showed a weak correlation with temperature 
($r = -0.119$), suggesting that density-dependent regulation and permanent aquatic habitats (water tanks) outweigh short-term climatic forcing. 
Conversely, in Dire Dawa, a strong negative correlation was observed between adult counts and precipitation ($r = -0.721, p < 0.001$). This 
potentially reflects a "wash-out" effect or a shift in habitat availability during high-rainfall events, which has been documented for urban 
container-breeding mosquitoes elsewhere \citep{2016.ST.OTI}. These results underscore the importance of local hydro-climatic contexts when 
modeling \textit{An. stephensi} stability.
\par
The spatial suitability maps generated via KDE correctly highlighted urban cores and high-building-density areas as the most suitable habitats. 
Validation against independent occurrence points yielded a Continuous Boyce Index (CBI) of \num{0.633} for Oromia, indicating moderate to high 
predictive performance. While the CBI for Dire Dawa was lower (\num{0.379}), this is likely an artifact of extreme data sparsity, as only one 
occurrence record remained within the study area after spatial clipping. The fact that \SI{100}{\percent} of the available points in Dire Dawa 
and \SI{14.3}{\percent} in Oromia fell into the top \SI{20}{\percent} of predicted suitability areas suggests that the model effectively 
captures high-risk zones, even with limited validation data.

\par 
The performance analysis (Table \ref{tab:performance}) demonstrates that the simulation framework is computationally efficient and can 
run on modest hardware. Both six-month simulations completed in less than an hour on a standard laptop (Intel Core i5, \SI{8}{\giga\byte} RAM), 
with average tick times below \SI{125}{\milli\second}. Peak memory usage remained within \SI{1.6}{\giga\byte}, well under the \SI{2}{\giga\byte} 
heap limit. These characteristics are critical for deployment in sub-Saharan African public health institutions, where access to high-performance 
computing clusters is often limited. The framework's ability to handle large geographical extents (e.g., the Oromia region with \num{60} 
million spatial cells) while maintaining acceptable performance is particularly noteworthy. By leveraging a lazy spatial grid and 
parallel agent processing, the model scales to region-wide simulations without requiring specialized hardware. 
\par
This opens the door for routine use by national malaria control programs and regional health observatories, enabling them to run scenario 
analyses, evaluate intervention strategies, and generate risk maps using only standard office computers. Future work will focus on reducing 
the memory footprint further and providing user-friendly interfaces to lower the technical barrier for public health officers in 
endemic areas.
\par
Several limitations should be acknowledged. The current movement rules utilize a random walk with a fixed step size, which may not fully capture 
complex host-seeking behaviors toward specific resource gradients. Furthermore, the model assumes human population density is the sole proxy 
for blood-meal availability. Given that \textit{An. stephensi} exhibits plastic feeding behavior, including zoophily, a multi-host formulation 
might improve fecundity estimates in peri-urban areas. 
\par
Despite these limitations, we believe the framework provides a valuable tool for designing targeted vector control interventions. The YAML-driven 
parameterisation allows for rapid adaptation to other regions or species. Future research should extend these simulations to multi-year runs 
to evaluate the impact of long-term climate variability and the effectiveness of larval source management, such as the systematic covering of 
water tanks, on population suppression.

\section{Conclusion}
\label{sec:conclusion}
This study presented an agent-based modelling framework for simulating the population dynamics of mosquitoes in low resource settings,
with a case study on \textit{Anopheles stephensi} in two Ethiopian regions (Dire Dawa & Oromia). The model successfully integrated 
temperature-dependent life-history traits with spatially explicit environmental data and a stage-structured Metzler matrix. The framework was applied 
to Dire Dawa and Oromia using a six-month simulation driven by ERA5-Land data. 
\par
The results demonstrated that the simulated populations reached steady or growing states consistent with field observations of established 
vectors. The model identified significant differences in regional growth factors ($\lambda_{\max}$ \numrange{1.003}{1.074}) and climatic 
sensitivities, highlighting the buffered nature of urban populations against environmental variability. Validation against independent 
occurrence data yielded Continuous Boyce Indices ranging from \numrange{0.379}{0.633}, confirming the model's ability to identify urban 
hotspots even in data-poor environments.
\par
While limitations regarding adult dispersal and host-seeking behavior remain, the framework provides a generic, reusable, and easily 
parameterisable tool for public health decision-making. By linking life-cycle biology with spatial infrastructure data, the model could
support the evidence-based planning of larval source management and urban malaria control programmes in recently invaded territories.


\section*{Data availability}
The code of the simulator and the data used for the case study are all available in a github directory:
\url{https://github.com/Vansnoden/monadvsim}.


\section*{Acknowledgments}
\par
This research was supported by the International Center of Insect Physiology and Ecology (ICIPE), 
the University of KwaZulu-Natal, and the International Institute of Tropical Agriculture (IITA). 
We acknowledge the providers of the open data used in this study, including GBIF, the Copernicus 
Climate Change Service (ERA5-Land), the Malaria Atlas Project, and the Vector Atlas consortium. 
We particularly thank the data contributors and curators who make their work publicly available, 
enabling research of this nature.


\bibliographystyle{apalike}
\bibliography{bibliography.bib}

\end{document}
